\documentclass[11pt,a4paper]{article}
\usepackage[utf8]{inputenc}
\usepackage[T1]{fontenc}
\usepackage[catalan]{babel}
\usepackage{lmodern}
\usepackage[sfdefault,lining]{FiraSans}
\usepackage[a4paper,top=2.3cm,bottom=2.2cm,left=1.8cm,right=1.8cm,headheight=1.3cm,headsep=0.4cm,footskip=1cm]{geometry}
\usepackage{amsmath,amssymb}
\usepackage{icomma}
\usepackage{graphicx}
\usepackage[table]{xcolor}
\usepackage{tikz}
\usepackage{fontawesome5}
\usepackage{ccicons}
\usepackage{fancyhdr}
\usepackage{lastpage}
\usepackage{tasks}
\usepackage{needspace}
\usepackage[most]{tcolorbox}
\usepackage{array}
\usepackage[hidelinks,pdftitle={Fitxa d'exercicis de la unitat 3: Proporcionalitat i percentatges. 3r d'ESO},pdfsubject={Matemàtiques 3r ESO},pdfkeywords={percentatges, augments, disminucions, interés compost, repartiments, mescles, moviments, proporcionalitat composta}]{hyperref}

\definecolor{navy}{HTML}{1B2748}
\definecolor{accent}{HTML}{E8603A}
\definecolor{teal}{HTML}{1F9E90}
\definecolor{soft}{HTML}{EEF1F8}
\definecolor{muted}{HTML}{5D6685}

\setlength{\parindent}{0pt}
\setlength{\parskip}{3pt}
\settasks{label=\alph*), label-format=\bfseries\color{navy}, label-width=1.3em, item-indent=1.8em, column-sep=0.6em, after-item-skip=2pt, before-skip=3pt, after-skip=4pt}

% ---------- capçalera i peu ----------
\newcommand{\llicencia}{\href{https://creativecommons.org/licenses/by-nc-sa/4.0/deed.ca}{\ccbyncsa}\ \ Llicència Creative Commons Reconeixement-NoComercial-CompartirIgual 4.0 Internacional (CC BY-NC-SA 4.0)}
\pagestyle{fancy}
\fancyhf{}
\fancyhead[L]{\raisebox{-0.35\height}{\includegraphics[height=0.95cm]{logo-gva.png}}}
\fancyhead[C]{\small\color{muted}Fitxa d'exercicis · Unitat 3: Proporcionalitat i percentatges · 3r d'ESO}
\fancyhead[R]{\raisebox{-0.35\height}{\includegraphics[height=0.95cm]{logo-vermellar.png}}}
\fancyfoot[L]{\scriptsize\color{muted}\llicencia}
\fancyfoot[R]{\small\color{muted}\thepage}
\renewcommand{\headrulewidth}{0.4pt}
\renewcommand{\headrule}{\color{muted!40}\hrule height \headrulewidth}
\fancypagestyle{primera}{\fancyhf{}\renewcommand{\headrulewidth}{0pt}
  \fancyfoot[L]{\scriptsize\color{muted}\llicencia}
  \fancyfoot[R]{\small\color{muted}\thepage}}

% ---------- nivells i exercicis ----------
\newcommand{\estrelles}[1]{{\color{accent}\ifcase#1\or\faStar\faStar[regular]\faStar[regular]\or\faStar\faStar\faStar[regular]\or\faStar\faStar\faStar\fi}}
\newcounter{exer}
\newcommand{\ex}[1]{\par\needspace{4\baselineskip}\medskip\stepcounter{exer}\noindent{\color{navy}\bfseries\theexer.}\ \ {\small\estrelles{#1}}\ \ }

\newcommand{\seccio}[2]{\par\needspace{10\baselineskip}\bigskip
  \begin{tcolorbox}[colback=navy,colframe=navy,coltext=white,arc=2mm,boxrule=0pt,left=3mm,right=3mm,top=1.2mm,bottom=1.2mm]
  {\large\bfseries #1\quad #2}
  \end{tcolorbox}\vspace{-2pt}}

% ---------- recta real ----------
\newcommand{\rectabuida}{\begin{tikzpicture}[x=0.42cm,baseline=0]
  \draw[-latex,muted] (-6.7,0)--(6.8,0);
  \foreach \i in {-6,...,6}{\draw[muted] (\i,0.12)--(\i,-0.12) node[below=-1pt,font=\scriptsize,text=muted]{$\i$};}
  \end{tikzpicture}}
% \rectaint{x1}{x2}{estil esquerra}{estil dreta}: estil = ple, buit, inf
\newcommand{\extrem}[2]{\ifx#2p\fill[accent] (#1,0) circle (0.13);\else\ifx#2b\filldraw[fill=white,draw=accent,line width=1.2pt] (#1,0) circle (0.13);\fi\fi}
\newcommand{\rectaint}[4]{\begin{tikzpicture}[x=0.42cm,baseline=0]
  \draw[-latex,muted] (-6.7,0)--(6.8,0);
  \foreach \i in {-6,...,6}{\draw[muted] (\i,0.12)--(\i,-0.12) node[below=-1pt,font=\scriptsize,text=muted]{$\i$};}
  \draw[accent,line width=3pt] (#1,0)--(#2,0);
  \ifx#3i\fill[accent] (#1-0.25,0)--(#1+0.15,0.2)--(#1+0.15,-0.2)--cycle;\else\extrem{#1}{#3}\fi
  \ifx#4i\fill[accent] (#2+0.25,0)--(#2-0.15,0.2)--(#2-0.15,-0.2)--cycle;\else\extrem{#2}{#4}\fi
  \end{tikzpicture}}

\newcolumntype{C}[1]{>{\centering\arraybackslash}p{#1}}

% =====================================================================
% Capçalera FentMates de les fitxes: un color per a cada curs.
% Els colors són els de la pàgina inicial (LEVELS a assets/js/materials.js);
% el logo ve de _eines/logo/pdf/ (python3 _eines/logo/genera.py).
%
% Ús, al final del preàmbul de la unitat:
%   % =====================================================================
% Capçalera FentMates de les fitxes: un color per a cada curs.
% Els colors són els de la pàgina inicial (LEVELS a assets/js/materials.js);
% el logo ve de _eines/logo/pdf/ (python3 _eines/logo/genera.py).
%
% Ús, al final del preàmbul de la unitat:
%   % =====================================================================
% Capçalera FentMates de les fitxes: un color per a cada curs.
% Els colors són els de la pàgina inicial (LEVELS a assets/js/materials.js);
% el logo ve de _eines/logo/pdf/ (python3 _eines/logo/genera.py).
%
% Ús, al final del preàmbul de la unitat:
%   \input{../comu/fentmates.tex}
%   \fmcurs{1eso}
% i, al començament del document:
%   \capcalera{Matemàtiques}{1r ESO}{Fitxa d'exercicis · Unitat 2}{Divisibilitat}
% Opcional abans de \input: \def\fmneed{14} (espai mínim abans de cada secció).
% =====================================================================
\makeatletter
\def\fm@logodir{../../logo/pdf}
\providecommand{\fmneed}{10}
\def\fm@cursnom{}

% clau del curs, color i color del text damunt del color
\newcommand{\fm@curs}[3]{\@namedef{fm@color@#1}{#2}\@namedef{fm@text@#1}{#3}}
\fm@curs{1eso}{E6B84A}{navy}
\fm@curs{2eso}{3B82F6}{white}
\fm@curs{3eso}{FF7A50}{white}
\fm@curs{4eso-a}{A78BFA}{white}
\fm@curs{4eso-b}{6D28D9}{white}
\fm@curs{1batx-mi}{9B1C1C}{white}
\fm@curs{1batx-ccss}{EF6B6B}{white}
\fm@curs{2batx-mii}{1E6B47}{white}
\fm@curs{2batx-ccss}{2BB3A3}{white}

\newcommand{\fmcurs}[1]{%
  \@ifundefined{fm@color@#1}{\PackageError{fentmates}{Curs desconegut: #1}{}}{}%
  \edef\fm@tmp{\noexpand\definecolor{curs}{HTML}{\@nameuse{fm@color@#1}}%
               \noexpand\colorlet{curstext}{\@nameuse{fm@text@#1}}}\fm@tmp
  \colorlet{cursfosc}{curs!72!black}%
  \def\fm@logo{\fm@logodir/vertical-#1.pdf}%
  \def\fm@marca{\fm@logodir/marca-#1.pdf}}

% ---------- capçalera de la primera pàgina ----------
% #1 matèria (petit), #2 curs, #3 tipus de document i unitat, #4 títol
\newcommand{\capcalera}[4]{%
  \gdef\fm@cursnom{#2}%
  \par\noindent
  \begin{tcolorbox}[colback=curs,colframe=curs,boxrule=0pt,arc=2.2mm,width=3.5cm,height=2.15cm,
      valign=center,left=3mm,right=2mm,top=1mm,bottom=1mm,nobeforeafter,box align=center]
    \color{curstext}{\footnotesize #1}\par\vspace{1pt}{\fontsize{19}{21}\selectfont\bfseries #2\par}%
  \end{tcolorbox}\hspace{5mm}%
  \begin{minipage}[c]{\dimexpr\linewidth-3.5cm-5mm-2.9cm\relax}
    {\color{cursfosc}\small\bfseries\MakeUppercase{#3}}\par\vspace{2pt}
    {\color{navy}\fontsize{20}{23}\selectfont\bfseries #4\par}
  \end{minipage}\hfill
  \begin{minipage}[c]{2.6cm}\raggedleft\includegraphics[width=2.6cm]{\fm@logo}\end{minipage}\par
  \vspace{2pt}}

% ---------- barres de secció del color del curs ----------
\renewcommand{\seccio}[2]{\par\needspace{\fmneed\baselineskip}\bigskip
  \begin{tcolorbox}[colback=curs,colframe=curs,coltext=curstext,arc=2mm,boxrule=0pt,left=3mm,right=3mm,top=1.2mm,bottom=1.2mm]
  {\large\bfseries #1\quad #2}
  \end{tcolorbox}\vspace{-2pt}}

% ---------- capçalera i peu de les pàgines ----------
% primera pàgina: logos institucionals al peu
\fancypagestyle{primera}{\fancyhf{}\renewcommand{\headrulewidth}{0pt}
  \fancyfoot[L]{\raisebox{-0.32\height}{\includegraphics[height=0.78cm]{logo-gva.png}}\hspace{5mm}%
    \raisebox{-0.32\height}{\includegraphics[height=0.78cm]{logo-vermellar.png}}}
  \fancyfoot[R]{\scriptsize\color{muted}\href{https://creativecommons.org/licenses/by-nc-sa/4.0/deed.ca}{\ccbyncsa}\ \ CC BY-NC-SA 4.0 · fentmates.org\qquad\small\thepage}}
% la resta: el símbol a l'esquerra i el curs a la dreta, amb el títol al centre com abans
\fancyhead[L]{\raisebox{-0.3\height}{\includegraphics[height=0.62cm]{\fm@marca}}}
\fancyhead[R]{\small\bfseries\color{cursfosc}\fm@cursnom}
\renewcommand{\headrule}{\color{curs!45}\hrule height \headrulewidth}
\makeatother

%   \fmcurs{1eso}
% i, al començament del document:
%   \capcalera{Matemàtiques}{1r ESO}{Fitxa d'exercicis · Unitat 2}{Divisibilitat}
% Opcional abans de \input: \def\fmneed{14} (espai mínim abans de cada secció).
% =====================================================================
\makeatletter
\def\fm@logodir{../../logo/pdf}
\providecommand{\fmneed}{10}
\def\fm@cursnom{}

% clau del curs, color i color del text damunt del color
\newcommand{\fm@curs}[3]{\@namedef{fm@color@#1}{#2}\@namedef{fm@text@#1}{#3}}
\fm@curs{1eso}{E6B84A}{navy}
\fm@curs{2eso}{3B82F6}{white}
\fm@curs{3eso}{FF7A50}{white}
\fm@curs{4eso-a}{A78BFA}{white}
\fm@curs{4eso-b}{6D28D9}{white}
\fm@curs{1batx-mi}{9B1C1C}{white}
\fm@curs{1batx-ccss}{EF6B6B}{white}
\fm@curs{2batx-mii}{1E6B47}{white}
\fm@curs{2batx-ccss}{2BB3A3}{white}

\newcommand{\fmcurs}[1]{%
  \@ifundefined{fm@color@#1}{\PackageError{fentmates}{Curs desconegut: #1}{}}{}%
  \edef\fm@tmp{\noexpand\definecolor{curs}{HTML}{\@nameuse{fm@color@#1}}%
               \noexpand\colorlet{curstext}{\@nameuse{fm@text@#1}}}\fm@tmp
  \colorlet{cursfosc}{curs!72!black}%
  \def\fm@logo{\fm@logodir/vertical-#1.pdf}%
  \def\fm@marca{\fm@logodir/marca-#1.pdf}}

% ---------- capçalera de la primera pàgina ----------
% #1 matèria (petit), #2 curs, #3 tipus de document i unitat, #4 títol
\newcommand{\capcalera}[4]{%
  \gdef\fm@cursnom{#2}%
  \par\noindent
  \begin{tcolorbox}[colback=curs,colframe=curs,boxrule=0pt,arc=2.2mm,width=3.5cm,height=2.15cm,
      valign=center,left=3mm,right=2mm,top=1mm,bottom=1mm,nobeforeafter,box align=center]
    \color{curstext}{\footnotesize #1}\par\vspace{1pt}{\fontsize{19}{21}\selectfont\bfseries #2\par}%
  \end{tcolorbox}\hspace{5mm}%
  \begin{minipage}[c]{\dimexpr\linewidth-3.5cm-5mm-2.9cm\relax}
    {\color{cursfosc}\small\bfseries\MakeUppercase{#3}}\par\vspace{2pt}
    {\color{navy}\fontsize{20}{23}\selectfont\bfseries #4\par}
  \end{minipage}\hfill
  \begin{minipage}[c]{2.6cm}\raggedleft\includegraphics[width=2.6cm]{\fm@logo}\end{minipage}\par
  \vspace{2pt}}

% ---------- barres de secció del color del curs ----------
\renewcommand{\seccio}[2]{\par\needspace{\fmneed\baselineskip}\bigskip
  \begin{tcolorbox}[colback=curs,colframe=curs,coltext=curstext,arc=2mm,boxrule=0pt,left=3mm,right=3mm,top=1.2mm,bottom=1.2mm]
  {\large\bfseries #1\quad #2}
  \end{tcolorbox}\vspace{-2pt}}

% ---------- capçalera i peu de les pàgines ----------
% primera pàgina: logos institucionals al peu
\fancypagestyle{primera}{\fancyhf{}\renewcommand{\headrulewidth}{0pt}
  \fancyfoot[L]{\raisebox{-0.32\height}{\includegraphics[height=0.78cm]{logo-gva.png}}\hspace{5mm}%
    \raisebox{-0.32\height}{\includegraphics[height=0.78cm]{logo-vermellar.png}}}
  \fancyfoot[R]{\scriptsize\color{muted}\href{https://creativecommons.org/licenses/by-nc-sa/4.0/deed.ca}{\ccbyncsa}\ \ CC BY-NC-SA 4.0 · fentmates.org\qquad\small\thepage}}
% la resta: el símbol a l'esquerra i el curs a la dreta, amb el títol al centre com abans
\fancyhead[L]{\raisebox{-0.3\height}{\includegraphics[height=0.62cm]{\fm@marca}}}
\fancyhead[R]{\small\bfseries\color{cursfosc}\fm@cursnom}
\renewcommand{\headrule}{\color{curs!45}\hrule height \headrulewidth}
\makeatother

%   \fmcurs{1eso}
% i, al començament del document:
%   \capcalera{Matemàtiques}{1r ESO}{Fitxa d'exercicis · Unitat 2}{Divisibilitat}
% Opcional abans de \input: \def\fmneed{14} (espai mínim abans de cada secció).
% =====================================================================
\makeatletter
\def\fm@logodir{../../logo/pdf}
\providecommand{\fmneed}{10}
\def\fm@cursnom{}

% clau del curs, color i color del text damunt del color
\newcommand{\fm@curs}[3]{\@namedef{fm@color@#1}{#2}\@namedef{fm@text@#1}{#3}}
\fm@curs{1eso}{E6B84A}{navy}
\fm@curs{2eso}{3B82F6}{white}
\fm@curs{3eso}{FF7A50}{white}
\fm@curs{4eso-a}{A78BFA}{white}
\fm@curs{4eso-b}{6D28D9}{white}
\fm@curs{1batx-mi}{9B1C1C}{white}
\fm@curs{1batx-ccss}{EF6B6B}{white}
\fm@curs{2batx-mii}{1E6B47}{white}
\fm@curs{2batx-ccss}{2BB3A3}{white}

\newcommand{\fmcurs}[1]{%
  \@ifundefined{fm@color@#1}{\PackageError{fentmates}{Curs desconegut: #1}{}}{}%
  \edef\fm@tmp{\noexpand\definecolor{curs}{HTML}{\@nameuse{fm@color@#1}}%
               \noexpand\colorlet{curstext}{\@nameuse{fm@text@#1}}}\fm@tmp
  \colorlet{cursfosc}{curs!72!black}%
  \def\fm@logo{\fm@logodir/vertical-#1.pdf}%
  \def\fm@marca{\fm@logodir/marca-#1.pdf}}

% ---------- capçalera de la primera pàgina ----------
% #1 matèria (petit), #2 curs, #3 tipus de document i unitat, #4 títol
\newcommand{\capcalera}[4]{%
  \gdef\fm@cursnom{#2}%
  \par\noindent
  \begin{tcolorbox}[colback=curs,colframe=curs,boxrule=0pt,arc=2.2mm,width=3.5cm,height=2.15cm,
      valign=center,left=3mm,right=2mm,top=1mm,bottom=1mm,nobeforeafter,box align=center]
    \color{curstext}{\footnotesize #1}\par\vspace{1pt}{\fontsize{19}{21}\selectfont\bfseries #2\par}%
  \end{tcolorbox}\hspace{5mm}%
  \begin{minipage}[c]{\dimexpr\linewidth-3.5cm-5mm-2.9cm\relax}
    {\color{cursfosc}\small\bfseries\MakeUppercase{#3}}\par\vspace{2pt}
    {\color{navy}\fontsize{20}{23}\selectfont\bfseries #4\par}
  \end{minipage}\hfill
  \begin{minipage}[c]{2.6cm}\raggedleft\includegraphics[width=2.6cm]{\fm@logo}\end{minipage}\par
  \vspace{2pt}}

% ---------- barres de secció del color del curs ----------
\renewcommand{\seccio}[2]{\par\needspace{\fmneed\baselineskip}\bigskip
  \begin{tcolorbox}[colback=curs,colframe=curs,coltext=curstext,arc=2mm,boxrule=0pt,left=3mm,right=3mm,top=1.2mm,bottom=1.2mm]
  {\large\bfseries #1\quad #2}
  \end{tcolorbox}\vspace{-2pt}}

% ---------- capçalera i peu de les pàgines ----------
% primera pàgina: logos institucionals al peu
\fancypagestyle{primera}{\fancyhf{}\renewcommand{\headrulewidth}{0pt}
  \fancyfoot[L]{\raisebox{-0.32\height}{\includegraphics[height=0.78cm]{logo-gva.png}}\hspace{5mm}%
    \raisebox{-0.32\height}{\includegraphics[height=0.78cm]{logo-vermellar.png}}}
  \fancyfoot[R]{\scriptsize\color{muted}\href{https://creativecommons.org/licenses/by-nc-sa/4.0/deed.ca}{\ccbyncsa}\ \ CC BY-NC-SA 4.0 · fentmates.org\qquad\small\thepage}}
% la resta: el símbol a l'esquerra i el curs a la dreta, amb el títol al centre com abans
\fancyhead[L]{\raisebox{-0.3\height}{\includegraphics[height=0.62cm]{\fm@marca}}}
\fancyhead[R]{\small\bfseries\color{cursfosc}\fm@cursnom}
\renewcommand{\headrule}{\color{curs!45}\hrule height \headrulewidth}
\makeatother

\fmcurs{3eso}

\begin{document}
\thispagestyle{primera}

% ---------- encapçalament ----------
\capcalera{Matemàtiques}{3r ESO}{Fitxa d'exercicis · Unitat 3}{Proporcionalitat i percentatges}
\vspace{10pt}
Nom i cognoms: \rule{8.4cm}{0.4pt}\hfill Grup: \rule{1.6cm}{0.4pt}\hfill Data: \rule{2.2cm}{0.4pt}
\vspace{8pt}

\begin{tcolorbox}[colback=soft,colframe=soft,arc=2mm,left=3mm,right=3mm,top=1.5mm,bottom=1.5mm]
\small Exercicis graduats per nivells de dificultat:\quad
\estrelles{1}\ bàsic \qquad \estrelles{2}\ mitjà \qquad \estrelles{3}\ avançat.\\[2pt]
Resol-los al quadern mostrant tots els passos. Les taules es poden completar directament en la fitxa. Pots fer servir la calculadora.
\end{tcolorbox}


% =====================================================
\seccio{1}{Percentatges, fraccions i decimals}

\ex{1} Completa la taula:\par\smallskip
{\renewcommand{\arraystretch}{1.75}\centering
\begin{tabular}{|C{4.2cm}|C{4.2cm}|C{4.2cm}|}
\hline
\rowcolor{soft}\textbf{Percentatge} & \textbf{Fracció irreductible} & \textbf{Decimal} \\ \hline
$40\,\%$ & & \\ \hline
 & $\dfrac34$ & \\ \hline
 & & $0,15$ \\ \hline
$8\,\%$ & & \\ \hline
 & $\dfrac{7}{20}$ & \\ \hline
 & & $1,25$ \\ \hline
$2,5\,\%$ & & \\ \hline
\end{tabular}\par}

\ex{1} Expressa com a percentatge:
\begin{tasks}(6)
  \task $0,32$
  \task $0,07$
  \task $1,5$
  \task $\dfrac35$
  \task $\dfrac{9}{25}$
  \task $\dfrac18$
\end{tasks}

\ex{2} Escriu com a fracció irreductible:
\begin{tasks}(6)
  \task $12\,\%$
  \task $65\,\%$
  \task $0,8\,\%$
  \task $12,5\,\%$
  \task $140\,\%$
  \task $0,25\,\%$
\end{tasks}

\ex{2} Ordena de menor a major:
\[
\frac25\qquad 0,38\qquad 41\,\%\qquad \frac{7}{16}\qquad 0,\overline{4}
\]

\ex{3} En el grup A han aprovat 18 dels 24 alumnes; en el B, 21 de 30, i en el C, 16 de 20. Quin percentatge d'aprovats té cada grup? Ordena els grups de millor a pitjor resultat.

% =====================================================
\seccio{2}{Càlcul de percentatges}

\ex{1} Calcula:
\begin{tasks}(3)
  \task El $20\,\%$ de $150$
  \task El $75\,\%$ de $64$
  \task El $5\,\%$ de $360$
  \task El $12\,\%$ de $250$
  \task El $150\,\%$ de $40$
  \task El $0,5\,\%$ de $2\,000$
\end{tasks}

\ex{1} Calcula mentalment:
\begin{tasks}(3)
  \task El $10\,\%$ de $470$
  \task El $50\,\%$ de $86$
  \task El $25\,\%$ de $120$
  \task L'$1\,\%$ de $3\,500$
  \task El $20\,\%$ de $45$
  \task El $15\,\%$ de $80$
\end{tasks}

\ex{2} En un institut de $850$ alumnes, el $46\,\%$ són xiques. Quantes xiques i quants xics hi ha?

\ex{2} Un jersei de $45$ € té un descompte del $30\,\%$. Quants euros ens descompten? Quant paguem?

\ex{2} S'ha preguntat a $1\,200$ persones pel seu esport preferit: el $35\,\%$ ha triat el futbol, el $28\,\%$ el bàsquet, el $15\,\%$ el tennis i la resta, altres esports. Quantes persones han triat cada opció?

\ex{3} Una empresa té un pressupost de $2\,500$ €. En gasta el $40\,\%$ en ordinadors i el $30\,\%$ del que li queda en programes. La resta es dedica a formació. Quants euros es destinen a cada cosa? Quin percentatge del pressupost total és la formació?

\ex{3} Calcula el $20\,\%$ del $30\,\%$ de $500$. És el mateix que el $30\,\%$ del $20\,\%$ de $500$? A quin percentatge únic de $500$ equival?

% =====================================================
\seccio{3}{Tant per cent que representa un nombre respecte d'un altre}

\ex{1} Quin percentatge representa:
\begin{tasks}(3)
  \task $12$ respecte de $48$
  \task $36$ respecte de $120$
  \task $9$ respecte de $12$
  \task $45$ respecte de $30$
  \task $3$ respecte de $200$
  \task $7$ respecte de $35$
\end{tasks}

\ex{1} En una classe de $25$ alumnes n'han aprovat $21$. Quin percentatge ha aprovat? I quin ha suspés?

\ex{2} Calcula el nombre en cada cas:
\begin{tasks}(2)
  \task El $15\,\%$ d'un nombre és $27$.
  \task El $40\,\%$ d'un nombre és $18$.
  \task El $120\,\%$ d'un nombre és $54$.
  \task El $2,5\,\%$ d'un nombre és $5$.
\end{tasks}

\ex{2} Un jugador ha marcat $18$ dels $24$ penals que ha llançat i un altre, $22$ de $30$. Qui té millor percentatge d'encert?

\ex{2} Per a reservar un viatge paguem una entrada de $36$ €, que és el $12\,\%$ del preu total. Quant costa el viatge? Quant ens queda per pagar?

\ex{3} Un embassament té una capacitat de $450$ hm$^3$ i ara en conté $162$ hm$^3$.
\begin{tasks}(1)
  \task A quin percentatge de la seua capacitat està?
  \task Després d'unes pluges arriba al $40\,\%$. Quants hm$^3$ d'aigua ha guanyat?
\end{tasks}

\ex{3} En una ciutat, $\frac38$ de la població són menors de $30$ anys i, d'aquests, el $40\,\%$ estudia. Quin percentatge de la població total són estudiants menors de $30$ anys?

% =====================================================
\seccio{4}{Augments i disminucions percentuals}

\ex{1} Completa la taula:\par\smallskip
{\renewcommand{\arraystretch}{1.6}\centering
\begin{tabular}{|C{3.2cm}|C{3.2cm}|C{3.2cm}|C{3.2cm}|}
\hline
\rowcolor{soft}\textbf{Quantitat inicial} & \textbf{Variació} & \textbf{Índex de variació} & \textbf{Quantitat final} \\ \hline
$200$ & augment del $15\,\%$ & & \\ \hline
$80$ & disminució del $25\,\%$ & & \\ \hline
$1\,200$ & augment del $3\,\%$ & & \\ \hline
$500$ & & $0,94$ & \\ \hline
$64$ & & & $80$ \\ \hline
\end{tabular}\par}

\ex{1} Quin augment o disminució percentual indica cada índex de variació?
\begin{tasks}(6)
  \task $1,12$
  \task $0,85$
  \task $1,03$
  \task $0,6$
  \task $2,5$
  \task $0,995$
\end{tasks}

\ex{2} Un ordinador costa $640$ € sense IVA. Quin és el preu final si l'IVA és del $21\,\%$?

\ex{2} El lloguer d'un pis és de $750$ € al mes i puja un $3,2\,\%$. Quin és el nou lloguer?

\ex{2} La població d'un poble ha passat de $12\,500$ a $11\,750$ habitants. Quin percentatge ha variat?

\ex{3} Una botiga apuja tots els preus un $10\,\%$ i una altra els rebaixa tots un $10\,\%$. Una mateixa jaqueta costava $90$ € en la primera botiga i $110$ € en la segona. En quina és més barata ara?

\ex{3} Un abonament al gimnàs costa $45$ € al mes. Si es paga l'any sencer d'una vegada, es fa un descompte del $15\,\%$ sobre el preu dels $12$ mesos. Quant es paga amb l'abonament anual? Quant s'estalvia?

% =====================================================
\seccio{5}{Càlcul de la quantitat inicial}

\ex{1} Calcula la quantitat inicial:
\begin{tasks}(2)
  \task Després d'augmentar un $20\,\%$ val $84$.
  \task Després de disminuir un $30\,\%$ val $56$.
  \task Després d'augmentar un $5\,\%$ val $420$.
  \task Després de disminuir un $12\,\%$ val $220$.
\end{tasks}

\ex{1} Amb el $21\,\%$ d'IVA inclòs, uns auriculars costen $181,50$ €. Quin és el preu sense IVA?

\ex{2} Rebaixada un $35\,\%$, una bicicleta costa $520$ €. Quin era el preu original? Quant ens estalviem?

\ex{2} Després d'una pujada del $3\,\%$, una treballadora cobra $1\,802,50$ €. Quin era el seu sou abans de la pujada?

\ex{2} Un cotxe d'ocasió es ven per $9\,600$ €, que és un $40\,\%$ menys del que valia nou. Quant valia nou?

\ex{3} Una factura de $96,80$ € inclou el $21\,\%$ d'IVA. Lluís diu: «L'IVA que hem pagat és el $21\,\%$ de $96,80$, és a dir, $20,33$ €».
\begin{tasks}(1)
  \task Calcula l'import sense IVA i l'IVA que s'ha pagat realment.
  \task Explica quin és l'error de Lluís.
\end{tasks}

\ex{3} Un dipòsit de $4\,000$ L de capacitat ha perdut el $8\,\%$ de l'aigua que tenia i ara en conté $3\,450$ L. Quants litres tenia abans de la pèrdua? A quin percentatge de la seua capacitat estava?

% =====================================================
\seccio{6}{Variacions percentuals encadenades}

\ex{1} Calcula l'índex de variació total i el percentatge de variació global:
\begin{tasks}(2)
  \task Un augment del $10\,\%$ i després un del $20\,\%$.
  \task Una disminució del $10\,\%$ i després una del $20\,\%$.
  \task Un augment del $50\,\%$ i després una disminució del $50\,\%$.
  \task Un augment del $25\,\%$ i després una disminució del $20\,\%$.
\end{tasks}

\ex{1} Un preu de $300$ € puja un $5\,\%$ i, més tard, un $8\,\%$. Quin és el preu final?

\ex{2} Un article de $120$ € està rebaixat un $30\,\%$ i a la caixa ens fan un $10\,\%$ de descompte addicional. Quant paguem? Quin és el descompte total que ens han fet?

\ex{2} El preu de la gasolina puja un $4\,\%$, després baixa un $6\,\%$ i, finalment, puja un $2\,\%$. Ha pujat o ha baixat en total? Quin percentatge?

\ex{2} Invertim $5\,000$ € en borsa. El primer any guanyen un $12\,\%$ i el segon perden un $15\,\%$. Quants diners tenim al final? Hem guanyat o perdut? Quin percentatge?

\ex{3} Un preu puja un $30\,\%$. Quin percentatge ha de baixar per a tornar al preu inicial?

\ex{3} Un preu puja un $20\,\%$ i després baixa un cert percentatge. El resultat és el mateix que si haguera baixat directament un $4\,\%$. Quin percentatge ha baixat la segona vegada?

% =====================================================
\seccio{7}{Interés compost}

\ex{1} Dipositem $3\,000$ € al $4\,\%$ d'interés compost anual durant $2$ anys. Quin capital tindrem? Quants interessos haurem guanyat?

\ex{1} Completa la taula d'un capital de $1\,000$ € al $5\,\%$ d'interés compost anual:\par\smallskip
{\renewcommand{\arraystretch}{1.6}\centering
\begin{tabular}{|C{2.2cm}|C{4cm}|C{4cm}|C{4cm}|}
\hline
\rowcolor{soft}\textbf{Any} & \textbf{Capital a l'inici} & \textbf{Interessos de l'any} & \textbf{Capital al final} \\ \hline
$1$ & $1\,000$ & & \\ \hline
$2$ & & & \\ \hline
$3$ & & & \\ \hline
\end{tabular}\par}

\ex{2} Calcula el capital final que s'obté en invertir $8\,000$ € al $2,5\,\%$ d'interés compost durant $5$ anys. Quants interessos s'obtindrien amb interés simple?

\ex{2} Quin capital hem de dipositar al $3\,\%$ d'interés compost anual per a tenir $10\,927,27$ € d'ací a $3$ anys?

\ex{2} Dipositem $12\,000$ € al $3\,\%$ anual amb abonament mensual dels interessos durant $2$ anys. Quin capital tindrem?

\ex{3} Quants anys han de passar perquè $5\,000$ € al $4\,\%$ d'interés compost anual superen els $6\,000$ €?

\ex{3} Volem invertir $10\,000$ € durant $4$ anys. El banc A ens ofereix un $3,5\,\%$ anual amb abonament anual i el banc B, un $3,4\,\%$ anual amb abonament mensual. Quina oferta és millor? Quina diferència hi ha?

% =====================================================
\seccio{8}{Repartiments proporcionals}

\ex{1} Reparteix $720$ de manera directament proporcional a $3$, $4$ i $5$.

\ex{1} Tres germans han treballat $10$, $15$ i $20$ hores en un negoci familiar i es reparteixen $1\,800$ € proporcionalment a les hores treballades. Quant rep cadascun?

\ex{2} Reparteix $1\,080$ de manera inversament proporcional a $3$, $4$ i $6$.

\ex{2} Un avi deixa $9\,000$ € als seus tres néts perquè se'ls repartisquen de manera inversament proporcional a les seues edats: $6$, $9$ i $18$ anys. Quant rep cadascun?

\ex{2} Tres socis han invertit $15\,000$ €, $20\,000$ € i $25\,000$ € en una empresa. Si els beneficis han sigut de $9\,600$ €, quant li correspon a cadascun?

\ex{3} Tres socis inverteixen en un negoci: A hi posa $6\,000$ € durant $8$ mesos; B, $4\,000$ € durant $12$ mesos, i C, $10\,000$ € durant $3$ mesos. Com s'han de repartir uns beneficis de $8\,400$ €?

\ex{3} Una prima de $4\,500$ € es reparteix entre dos treballadors de manera directament proporcional a les hores extres que han fet ($30$ i $40$) i inversament proporcional a les faltes que han tingut ($3$ i $2$). Quant rep cadascun?

\ex{3} Un premi es reparteix de manera directament proporcional a $2$, $3$ i $5$. Si el tercer rep $450$ € més que el primer, de quant és el premi? Quant rep cadascun?

% =====================================================
\seccio{9}{Problemes de mescles}

\ex{1} Es mesclen $5$ kg de cafè de $6$ €/kg amb $3$ kg de cafè de $10$ €/kg. A quin preu ix el quilo de mescla?

\ex{1} Es mesclen $12$ L de vi de $2,50$ €/L amb $8$ L de vi de $4$ €/L. Quin és el preu del litre de mescla?

\ex{2} Quants quilos de cafè de $9$ €/kg s'han de mesclar amb $20$ kg de cafè de $6$ €/kg per a obtindre una mescla de $7$ €/kg?

\ex{2} Es mesclen $3$ L d'un suc amb un $60\,\%$ de fruita i $2$ L d'una beguda amb un $20\,\%$ de fruita. Quin percentatge de fruita té la mescla?

\ex{2} Un comerciant mescla $30$ kg d'arròs de $1,80$ €/kg amb $20$ kg d'arròs de $2,40$ €/kg. A quin preu ha de vendre el quilo de mescla per a guanyar un $25\,\%$?

\ex{3} Es fonen $40$ g d'or de llei $0,750$ amb $10$ g d'or pur. Quina és la llei de l'aliatge obtingut? Quants quirats té?

\ex{3} Quants grams d'or pur s'han d'afegir a $60$ g d'or de $18$ quirats per a obtindre or de $20$ quirats?

\ex{3} Tenim $15$ L d'alcohol al $70\,\%$. Quanta aigua cal afegir-hi perquè quede al $40\,\%$?

% =====================================================
\seccio{10}{Problemes de moviments}

\ex{1} Expressa $54$ km/h en m/s i $25$ m/s en km/h.

\ex{1} Un tren circula a $120$ km/h.
\begin{tasks}(1)
  \task Quant de temps tarda a recórrer $330$ km? Expressa-ho en hores i minuts.
  \task Quina distància recorre en $1$ h $20$ min?
\end{tasks}

\ex{2} Dues ciutats estan a $360$ km. A les $9{:}00$ ix un cotxe de cada ciutat, l'un cap a l'altre, a $70$ km/h i a $90$ km/h. A quina hora es troben? A quina distància de la ciutat del primer cotxe?

\ex{2} Un vianant ix a $5$ km/h. Dues hores més tard, ix del mateix lloc i pel mateix camí un ciclista a $15$ km/h. Quant de temps tarda a atrapar-lo? A quina distància del punt de partida?

\ex{2} Dos ciclistes ixen alhora del mateix punt en sentits contraris, a $18$ km/h i a $22$ km/h. Quant de temps passarà fins que estiguen a $100$ km l'un de l'altre?

\ex{3} Les ciutats A i B estan a $500$ km. A les $7{:}30$ ix un camió d'A cap a B a $80$ km/h i a les $9{:}00$ ix un cotxe de B cap a A a $110$ km/h. A quina hora es creuen? A quina distància d'A?

\ex{3} Un tren de mercaderies ix d'una estació a $60$ km/h. Una hora i mitja més tard, ix de la mateixa estació, per una via paral·lela i en el mateix sentit, un tren ràpid a $150$ km/h. Quant tarda el ràpid a abastar el de mercaderies? A quants quilòmetres de l'estació?

\ex{3} Dues atletes ixen alhora i en el mateix sentit en una pista de $400$ m, a $6$ m/s i a $5$ m/s. Quant de temps tarda la més ràpida a doblar l'altra? Quantes voltes ha fet cadascuna en aquest temps?

% =====================================================
\seccio{11}{Proporcionalitat composta}

\ex{1} Indica si cada parella de magnituds és directament proporcional, inversament proporcional o no proporcional:
\begin{tasks}(2)
  \task La velocitat d'un cotxe i el temps que tarda a fer un trajecte.
  \task Els quilos de fruita que comprem i el preu que paguem.
  \task L'edat d'una persona i la seua altura.
  \task El nombre d'obrers i els dies que tarden a fer una obra.
  \task El costat d'un quadrat i la seua àrea.
  \task El nombre de grifs oberts i els litres que aboquen en una hora.
\end{tasks}

\ex{1} Resol:
\begin{tasks}(1)
  \task Si $4$ entrades de cine costen $26$ €, quant costen $7$ entrades?
  \task $8$ operaris fan un treball en $15$ dies. Quants dies tardarien $12$ operaris?
\end{tasks}

\ex{2} $6$ aixetes obertes durant $4$ hores aboquen $3\,600$ L d'aigua. Quants litres abocaran $9$ aixetes obertes durant $5$ hores?

\ex{2} $10$ vaques mengen $1\,800$ kg de farratge en $12$ dies. Quant de farratge menjaran $15$ vaques en $20$ dies?

\ex{2} $5$ obrers, treballant $8$ hores diàries, fan una obra en $18$ dies. Quants dies tardarien $6$ obrers treballant $10$ hores diàries?

\ex{3} $12$ impressores, funcionant $6$ hores diàries, imprimeixen $54\,000$ fulls en $5$ dies. Quants dies necessitaran $10$ impressores, funcionant $8$ hores diàries, per a imprimir $72\,000$ fulls?

\ex{3} Amb $120$ kg de pinso s'alimenten $30$ gallines durant $20$ dies. Quantes gallines es poden alimentar durant $25$ dies amb $90$ kg de pinso?

\ex{3} $8$ obrers, treballant $9$ hores diàries, construeixen $120$ m de mur en $15$ dies. Quants dies tardaran $10$ obrers, treballant $6$ hores diàries, a construir $200$ m de mur?

\vfill
\begin{center}\color{muted}\small\faCheckCircle[regular]\ \ Ànim! Si et costa algun exercici, repassa els apunts interactius de la unitat.\end{center}

\end{document}
